\section{Trabajos a futuro}

Este trabajo sienta las bases para múltiples líneas de investigación futuras que permitirán
profundizar y extender los resultados obtenidos. Los avances propuestos abarcan desde la
validación experimental del controlador diseñado hasta el desarrollo de estrategias de control
más sofisticadas que contemplen de forma explícita las no linealidades identificadas en el
sistema de levitación magnética ECP-730.

\begin{itemize}

    \item \textbf{Validación experimental en el dispositivo ECP-730.}
    Implementar físicamente el controlador PII diseñado en el dispositivo ECP-730 real,
    verificando que los parámetros de simulación $(a,\, b,\, c_1,\, m)$ corresponden al
    sistema físico. En caso de discrepancia, se deberá llevar a cabo un proceso de
    \emph{identificación experimental} que ajuste dichos parámetros mediante ensayos de
    entrada-salida en lazo abierto, consolidando así la correspondencia entre el modelo
    matemático y el comportamiento observado en el laboratorio.

    \item \textbf{Comparativa de modelos de fuerza electromagnética.}
    Evaluar los diferentes modelos de la \emph{fuerza de campo} $f_{\mathrm{fld}}$
    reportados en la literatura ---incluyendo los propuestos por Zhao \textit{et al.},
    El Hajjaji \textit{et al.}, Zhang \textit{et al.} y Al-Muthair \textit{et al.}---
    en términos de precisión de predicción y complejidad computacional para el SLM ECP-730.
    Dicha comparativa permitirá seleccionar el modelo que ofrezca el mejor equilibrio entre
    fidelidad física y viabilidad de implementación en tiempo real.

    \item \textbf{Comparativa de modelos de fricción.}
    Comparar el \emph{modelo de Karnopp} empleado en este trabajo con otros modelos
    dinámicos de fricción, tales como LuGre, Stribeck, Dahl y elasto-plástico, evaluando
    su fidelidad numérica y costo computacional en simulaciones de lazo cerrado. Este
    análisis proporcionará criterios objetivos para seleccionar el modelo más adecuado cuando
    se requiera reproducir con mayor exactitud los transitorios de ruptura estática
    (\textit{breakaway}) observados experimentalmente.

    \item \textbf{Identificación de parámetros de la zona muerta dinámica.}
    Los parámetros que caracterizan la zona muerta dinámica no centrada ---$m_r$, $m_l$,
    $b_r$, $b_l$, $h_r$ y $h_l$--- son desconocidos y se supusieron como hipótesis de
    modelado en este trabajo. Se propone desarrollar un \emph{algoritmo de identificación
    en línea}, basado, por ejemplo, en mínimos cuadrados recursivos o en esquemas
    adaptativos con modelo de referencia (MRAC), que estime dichos parámetros durante la operación normal del
    sistema. La disponibilidad de estimaciones precisas permitiría incorporar una
    compensación activa de la zona muerta directamente en la ley de control.

    \item \textbf{Estructuras de control frente al ciclo límite.}
    El barrido de familias PI y PII mostró que la amplitud del ciclo límite depende sobre todo
    de la ganancia proporcional y no del número de integradores. Se propone explorar
    estructuras que actúen directamente sobre el ciclo límite, como la compensación de la
    fricción estática, el control conmutado o el \textit{dither}, y evaluar si la ventaja del PII
    en rampa se vuelve relevante con referencias de pendiente mucho mayor o con una planta de
    menor fricción estática. Conviene además rediseñar el controlador en el centro del
    intervalo de operación ($x_1^*\approx-2{,}5$~cm), donde el margen de fase del diseño actual
    se reduce a $51^\circ$.

    \item \textbf{Control robusto mediante $H_\infty$ o QFT.}
    Diseñar un controlador robusto que garantice \emph{estabilidad y desempeño} ante
    variaciones paramétricas (incertidumbre en $b$, $c_1$ y $m$) y perturbaciones no
    modeladas, empleando técnicas de control $H_\infty$ o \textit{Quantitative Feedback
    Theory} (QFT). Dado que el punto de equilibrio del sistema es inherentemente inestable,
    cuantificar y acotar explícitamente los márgenes de robustez resulta de particular
    relevancia para garantizar la integridad del sistema ante condiciones de operación
    variables.

    \item \textbf{Control con ganancia programada para múltiples puntos de equilibrio.}
    Extender el diseño a un esquema de \emph{control con ganancia programada}
    (\textit{gain scheduling}) que ajuste automáticamente los parámetros del controlador
    PII conforme la posición de referencia varía dentro del rango operativo
    $x_1^* \in [-3{,}6;\,-0{,}5]$~cm, en el que los cambios de referencia son viables con $u\le 3{,}5$~V. Dado que el modelo linealizado ---y en consecuencia
    los polos del sistema--- cambia significativamente con el punto de operación, un
    esquema de este tipo permitiría mantener un desempeño uniforme en todo el rango operativo
    sin necesidad de rediseñar el controlador de manera manual.

    \item \textbf{Observador de estado para estimación de velocidad.}
    El PII propuesto solo requiere la posición, pero varias de las estrategias anteriores, como la compensación de la fricción o la linealización por retroalimentación, necesitan conocer la \emph{velocidad} $x_2$, que no se mide directamente: obtenerla derivando la posición amplifica el ruido de medición. Alcántara et al.~\cite{alcantara}
    abordaron este problema para la configuración repulsiva del ECP-730, diseñando e
    implementando tanto un observador lineal de Luenberger como un observador no lineal
    basado en la estructura del sistema, y evaluaron su desempeño en tiempo real.
    Se propone extender esa metodología a la configuración inestable, diseñando un observador
    que sea robusto ante la inestabilidad del sistema no observado, y evaluar cuantitativamente
    su impacto sobre el desempeño y la estabilidad del lazo cerrado con el controlador PII.

    \item \textbf{Control conmutado como estrategia complementaria.}
    El esquema de \emph{control conmutado} propuesto por Alcántara et al.~\cite{alcantara}
    ---que alterna entre un controlador con acción integral (para garantizar bajo error en estado
    estacionario) y uno sin integral (para suprimir el atascamiento-deslizamiento)--- demostró
    ser efectivo en la configuración repulsiva y constituye un punto de comparación relevante
    para el PII propuesto en este trabajo. Se propone implementar el controlador conmutado
    PI--PD, también analizado por Pérez-Gómez~\cite{perez} para el motor de CD, en la
    configuración inestable del SLM, a fin de realizar una comparativa cuantitativa entre
    ambas estrategias en términos de tiempo de atascamiento, error en estado estacionario
    y complejidad de implementación.

\end{itemize}
